% Präambel
\documentclass[11pt,			% Schriftgröße
a4paper,						% Papierformat
twoside, 						% einseitiges (oneside) oder zweiseitiges (twoside) Dokument
listof=totoc, 					% Tabellen- und Abbildungsverzeichnis ins Inhaltsverzeichnis
bibliography=totoc,				% Literaturverzeichnis ins Inhaltsverzeichnis aufnehmen
titlepage, 						% Titlepage-Umgebung statt \maketitle
headsepline, 					% horizontale Linie unter Kolumnentitel
%abstracton,					% Überschrift beim Abstract einschalten, Abstract muss dazu in {abstract}-Umgebung stehen
DIV18,							% auskommentieren, um den Seitenspiegel zu vergrößern
BCOR6mm,						% Bindekorrektur, die den Seitenspiegel um 6mm nach rechts verschiebt,
cleardoublepage=empty,			% Stil einer leeren eingefügten Seite bei Kapitelwechsel
parskip							% Absatzabstand bei Absatzwechsel einfügen
]{scrbook}			
\usepackage{ucs} 				% Dokument in utf8-Codierung schreiben und speichern
\usepackage[utf8x]{inputenc} 	% ermöglicht die direkte Eingabe von Umlauten
\usepackage[ngerman]{babel} 	% deutsche Trennungsregeln und Übersetzung der festcodierten Überschriften
\usepackage[T1]{fontenc} 		% Ausgabe aller zeichen in einer T1-Codierung (wichtig für die Ausgabe von Umlauten!)
\setlength{\parindent}{0ex} 	% bei neuem Abschnitt nicht einrücken
\linespread{1.2}\selectfont     % Zeilenabstand erhöhen - größere Werte als 1.2 nicht verwenden!!
\usepackage{scrpage2}			% SCR Headings verwenden
\setheadsepline{0.4pt}			% Kopfzeile Linien oben
\setfootsepline{0.4pt}			% Kopfzeile Linien unten
\pagestyle{scrheadings}			% SCR Headings einschalten
\usepackage{graphicx}  			% Einbinden von Grafiken erlauben
\usepackage[format=hang,		% Formatierungen von Unter- / Überschriften
font=normal,
labelfont=bf,
justification=RaggedRight,
singlelinecheck=true,
aboveskip=1mm
]{caption}


% Please add the following required packages to your document preamble:
\usepackage[table,xcdraw]{xcolor}
% If you use beamer only pass "xcolor=table" option, i.e. \documentclass[xcolor=table]{beamer}
% Note: It may be necessary to compile the document several times to get a multi-page table to line up properly
\usepackage{booktabs}
\usepackage{longtable}
\usepackage{pdfpages}



\usepackage{enumitem}			% Erlaubt Änderung der Nummerierung in der Umgebung enumerate

\usepackage{amsmath}			% Ergänzungen für Formeln
\usepackage{textcomp} 			% zum Einsatz von Eurozeichen u. a. Symbolen
\newcommand*\diff{\mathop{}\!\mathrm{d}}	% Differentialzeichen
\newcommand*\Diff[1]{\mathop{}\!\mathrm{d^#1}} % Differentialzeichen höherer Ableitung
\newcommand*\jj{\mathop{}\!\mathrm{j}}	% Komplexe Zahl j

\usepackage[					% Einstellunge Paket hyperref
hyperfootnotes=false			% im pfd-Output Fußnoten nicht verlinken
]{hyperref}

\usepackage{makeidx}			% Paket zur Erstellung eines Index
\usepackage[intoc]{nomencl} 	% zur Erstellung des Abkürzungsberzeichnisses

\usepackage[					% Einstellungen für Fußnoten
bottom,							% Ausrichtung unten
multiple,						% Trennung durch Seperator bei mehreren Fußnoten
hang,
marginal
]{footmisc}

\usepackage{calc}				% Paket zum Berechnen von Längen z.B. 0.8\linewidth

\usepackage{xcolor} 			% einfache Verwendung von Farben in nahezu allen Farbmodellen

\usepackage{listings}			% Darstellung von Quellcode mit den Umgebungen {lstlisting}, \lstinline und \lstinputlisting
\lstset{literate=				% Damit können Umlaute innerhalb Listings geschrieben werden
	{Ö}{{\"O}}1
	{Ä}{{\"A}}1
	{Ü}{{\"U}}1
	{ß}{{\ss}}1
	{ü}{{\"u}}1
	{ä}{{\"a}}1
	{ö}{{\"o}}1
}
\definecolor{mygreen}{rgb}{0,0.6,0}
\definecolor{mygray}{rgb}{0.5,0.5,0.5}
\definecolor{mymauve}{rgb}{0.58,0,0.82}
\lstset{ %
	backgroundcolor=\color{white},   % choose the background color; you must add \usepackage{color} or \usepackage{xcolor}; should come as last argument
	basicstyle=\footnotesize,        % the size of the fonts that are used for the code
	breakatwhitespace=false,         % sets if automatic breaks should only happen at whitespace
	breaklines=true,                 % sets automatic line breaking
	captionpos=t,                    % sets the caption-position to (b) bottom or (t) top
	commentstyle=\color{mygreen},    % comment style
	deletekeywords={...},            % if you want to delete keywords from the given language
	escapeinside={\%*}{*)},          % if you want to add LaTeX within your code
	escapeinside={(*@}{@*)},
	extendedchars=true,              % lets you use non-ASCII characters; for 8-bits encodings only, does not work with UTF-8
	frame=none,	                   	% "single" adds a frame around the code; "none"
	keepspaces=true,                 % keeps spaces in text, useful for keeping indentation of code (possibly needs columns=flexible)
	keywordstyle=\color{blue},       % keyword style
	language=[LaTeX]TeX,             % the language of the code
	morekeywords={*,nomenclature},   % if you want to add more keywords to the set
	numbers=left,                    % where to put the line-numbers; possible values are (none, left, right)
	numbersep=5pt,                   % how far the line-numbers are from the code
	numberstyle=\tiny\color{mygray}, % the style that is used for the line-numbers
	rulecolor=\color{black},         % if not set, the frame-color may be changed on line-breaks within not-black text (e.g. comments (green here))
	showspaces=false,                % show spaces everywhere adding particular underscores; it overrides 'showstringspaces'
	showstringspaces=false,          % underline spaces within strings only
	showtabs=false,                  % show tabs within strings adding particular underscores
	stepnumber=1,                    % the step between two line-numbers. If it's 1, each line will be numbered
	stringstyle=\color{mymauve},     % string literal style
	tabsize=2,	                   % sets default tabsize to 2 spaces
	title=\lstname                   % show the filename of files included with \lstinputlisting; also try caption instead of title
}

\makeindex						% Indexverzeichnis erstellen
\makenomenclature				% Abkürzungsverzeichnis erstellen

% -----------------------------------------------------------------------------------------------------------------
% Zum Aktualisieren des Abkürzungsverzeichnisses (Nomenklatur) bitte auf der Kommandozeile folgenden Befehl aufrufen :
% makeindex <Dateiname>.nlo -s nomencl.ist -o <Dateiname>.nls
% Oder besser: Kann in TexStudio unter Tools-Benutzer als Shortlink angelegt werden
% Konfiguration unter: Optionen-Erzeugen-Benutzerbefehle: makeindex -s nomencl.ist -t %.nlg -o %.nls %.nlo
% -----------------------------------------------------------------------------------------------------------------

% Hier die persönlichen Daten eingeben:

\newcommand{\titel}{BlindControl}
\newcommand{\untertitel}{Eine smarte Steuerung für Jalousien}
\newcommand{\arbeit}{Projektdokumentation}
\newcommand{\studiengang}{Elektrotechnik}
\newcommand{\studienrichtung}{Fahrzeugelektronik}
\newcommand{\studienschwerpunkt}{}
\newcommand{\autor}{Maximilian Bachmann, Marvin Eisenmann, Florian Vetter}
\newcommand{\matrikelnr}{123 456}
\newcommand{\kurs}{TFE17-2}
\newcommand{\firma}{Firma (Angabe entfällt ggf. bei Studienarbeit)}
\newcommand{\abgabe}{08.03.2019}
\newcommand{\betreuerdhbw}{Hans Juergen Herpel}
\newcommand{\betreuerfirma}{Gutachter der Firma (Angabe entfällt ggf. bei Studienarbeit)}
\newcommand{\jahr}{2019}			% für Angabe im Copyright-Vermerk der Titelseite

% Folgende Zeilen definieren Abkürzungen, um Befehle schneller eingeben zu können
\newcommand{\ua}{\mbox{u.\,a.\ }}
\newcommand{\zB}{\mbox{z.\,B.\ }}
\newcommand{\bs}{$\backslash$}

% Folgende Zeilen weden benötigt, um Tikz und PGF-Plot-Grafiken einzubinden
\usepackage{pgfplots}
\usepackage{pgfplotstable}
\pgfplotsset{compat=newest}
\usepgfplotslibrary{smithchart}
\usepackage{tikz}						% Tikz sollte nach Listings Pakete geladen werden.
\usetikzlibrary{arrows}

\hyphenation{Schrift-ar-ten}


% -------------------------------------------------------------------------------------------
%                     Beginn des Dokumenteninhalts
% -------------------------------------------------------------------------------------------
\begin{document}
\setcounter{secnumdepth}{3}				% Nummerierungstiefe fürs Inhaltsverzeichnis
\setcounter{tocdepth}{3}
\sffamily								% für die Titelei serifenlose Schrift verwenden

% ------------------------------ Titelei -----------------------------------------------------

\thispagestyle{plain}
\begin{titlepage}
\enlargethispage{4.0cm}
\sffamily 								% Serifenlose Grundschrift für die Titelseite einstellen

\parbox{0.5\linewidth}{
\begin{flushleft}
% Hier ggf. ein Logo der Firma
\end{flushleft}
}
\parbox{0.5\linewidth}{
\begin{flushright}
	\includegraphics[width=0.4\linewidth]{images/DHBW_d_R_FN_46mm_4c}\\[5ex]
\end{flushright}
}
				

\begin{center}

\huge{\textsc{\textbf{\titel}}}\\[1.5ex]
\Large{\textbf{\untertitel}}\\[5ex]
\LARGE{\textbf{\arbeit}}\\[2ex]
%\normalsize{für die Prüfung zum\\[1ex] Bachelor of Engineering}\\[3ex]
\Large{Studiengang \studiengang}\\[2ex]
\large{Studienrichtung \studienrichtung}\\[1ex]
%\large{ \studienschwerpunkt}\\[2ex]
\normalsize{Duale Hochschule Baden-Württemberg Ravensburg, Campus Friedrichshafen}\\[5ex]
von\\[1ex] \autor \\[20ex]


\end{center}

\begin{flushleft}

\begin{tabular}{ll}
Abgabedatum:					& \quad \abgabe \\
Bearbeitungszeitraum:		   		& \quad 07.01.2019 - 07.03.2019   \\ 
%Matrikelnummer: 			& \quad \matrikelnr \\ 
Kurs: 							& \quad \kurs \\
%Ausbildungsfirma:	 			& \quad \firma \\ 
%Betreuer der Ausbildungsfirma:  & \quad \betreuerfirma \\ 
%Gutachter der Dualen Hochschule: & \quad \betreuerdhbw \\ [5ex]

\end{tabular} 



\small
Copyrightvermerk:\\

Dieses Werk einschließlich seiner Teile ist \textbf{urheberrechtlich geschützt}. Jede Verwertung außerhalb der engen Grenzen des Urheberrechtgesetzes ist ohne Zustimmung des Autors unzulässig und strafbar. Das gilt insbesondere für Vervielfältigungen, Übersetzungen, Mikroverfilmungen sowie die Einspeicherung und Verarbeitung in elektronischen Systemen.
\end{flushleft}
\begin{flushright}
\copyright{} \jahr
\end{flushright}
\end{titlepage}

\cleardoublepage 				% erzeugt die Titelseite
\pagenumbering{roman}					% kleine, römische Seitenzahlen für Titelei
%%% Ggf. folgende Zeile auskommentieren, falls der Sperrvermerk gewünscht ist.
%\chapter*{Sperrvermerk} %*-Variante sorgt dafür, das der Sperrvermerk nicht im Inhaltsverzeichnis auftaucht
%gemäß Ziffer 1.1.13 der Anlage 1 zu §§ 3, 4 und 5  der Studien- und Prüfungsordnung für die Bachelorstudiengänge im Studienbereich Technik der Dualen Hochschule Baden-Württemberg vom 29.09.2017 in der Fassung vom 25.07.2018:
%
%Der Inhalt dieser Arbeit darf weder als Ganzes noch in Auszügen Personen außerhalb des Prüfungsprozesses und des Evaluationsverfahrens zugänglich gemacht werden, sofern keine anders lautende Genehmigung vom Dualen Partner vorliegt.
%
%Musterstadt, den \today \\[4ex]
%
%\rule[-0.2cm]{5cm}{0.5pt} \\
%
%\textsc{\autor} \\[10ex]

\chapter*{Erklärung} %*-Variante sorgt dafür, dass die Erklärung nicht im Inhaltsverzeichnis auftaucht

gemäß Ziffer 1.1.13 der Anlage 1 zu §§ 3, 4 und 5  der Studien- und Prüfungsordnung für die Bachelorstudiengänge im Studienbereich Technik der Dualen Hochschule Baden-Württemberg vom 29.09.2017 in der Fassung vom 25.07.2018.

Ich versichere hiermit, dass ich meine Bachelorarbeit (bzw. Projektarbeit oder Studienarbeit bzw. Hausarbeit) mit dem Thema:

\begin{quote}
	\textit{\titel} -\textit{ \untertitel }
\end{quote}

selbstständig verfasst und keine anderen als die angegebenen Quellen und Hilfsmittel benutzt habe. Ich versichere zudem, dass die eingereichte elektronische Fassung mit der gedruckten Fassung übereinstimmt.\\[10ex]

Musterstadt, den \today \\[4ex]

\rule[-0.2cm]{5cm}{0.5pt} \\

\textsc{\autor} \\[10ex]

\cleardoublepage 				% Einbinden der eidestattlichen Erklärung
%\chapter*{Kurzfassung} %*-Variante sorgt dafür, das Abstract nicht im Inhaltsverzeichnis auftaucht

Text Kurzfassung

\cleardoublepage
   			% Einbinden des Abstracts

\tableofcontents						% Erzeugen des Inhalsverzeichnisses
\cleardoublepage

% --------------------------------------------------------------------------------------------
%                    Inhalt der Bachelorarbeit
%---------------------------------------------------------------------------------------------
\pagenumbering{arabic}					% arabische Seitenzahlen für den Hauptteil

\rmfamily

\chapter{Einführung}
\label{cha:Einführung}
Das Projekt BlindControl entstand im Rahmen der Mikrocontroller Vorlesung des Dozenten Hans Jürgen Herpel an der DHBW Ravensburg Campus Friedrichshafen. Ziel ist es, eine Projektentwicklung anhand eines selbstgewählten Beispiels zu durchlaufen, um die entsprechenden Vorgänge zu erlernen und zu dokumentieren. Das fertige Projekt kann in dem zugehörigen \href{https://github.com/maxbachmann/BlindControl}{Git Repository} auf Github angesehen und verwendet werden. Zudem gibt es zu dem Projekt ein ausführliches \href{https://github.com/maxbachmann/BlindControl/wiki}{Github Wiki}.

% Was ist das Problem?
\section{Heranführung an die Problemstellung}
Die dem Projekt zugrundeliegende Idee ist die automatische Ansteuerung einer Jalousie mittels Home Assistant. Viele Jalousien sind nur manuell zu steuern oder nur mit teurer Technik automatisiert. Dennoch birgt eine automatische Steuerung einen hohen Comfort, auf den auch andere Anwender nicht verzichten sollten. Anwendungsbeispiele sind Folgende:
\begin{itemize}
	\item Bei zu hoher Sonneneinstrahlung verdunkelt die Jalousie automatisch, um eine angenehme Raumtemperatur auch im Sommer zu wahren.
	\item Jalousien sind durch ihren Aufbau sehr windanfällig und können leicht Schaden nehmen, oder gar anwesende Personen gefährden, sollte ein Sturm aufziehen. Auch hier kann eine Steuerung reagieren und die Jalousien automatisiert hochfahren und so zu schützen.
	\item Zusätzlich zu der automatisierten Steuerung kann es auch vorkommen, dass ein Nutzer beispielsweise auf dem Sofa liegt und durch die Sonnen geblendet wird. Um nicht aufstehen und zum Schalter an der Wand laufen zu müssen, könnte eine Heimautomatisierung zum Einsatz kommen, in der eine manuelle Steuerung der Jalousien integriert ist. Damit könnten diese mit dem Smartphone, das viele Benutzer immer bei sich haben, gesteuert werden.
	\item Auch die nachträgliche Installation solcher Technik kann Probleme bereiten, wenn beispielsweise Leitungen von der Jalousie zu dem Steuerungsserver gelegt werden müssen. Dagegen ist eine auf Funk basierte Kommunikation der einzelnen Komponenten optimal.
\end{itemize}
Auf Grundlage der genannten Anwendungssituationen wurden für das Projekt BlindControl die Anforderungen(Requirements) entworfen, sowie die Hardware und Software ausgewählt.

% Wie sieht der Lösungsraum aus?
\section{Lösungsansätze}
Um alle genannten Funktionen und Eigenschaften zusammen nutzen zu können, wird eine zentrale Steuereinheit benötigt, die alle Jalousien verwaltet. Auf diese sollte der Benutzer zugreifen können, um die manuelle Steuerung zu übernehmen. Zudem sollten hier alle Befehle an die Jalousien zentral gesendet werden, die damit eine eigene Ansteuerung brauchen. Um auf Umwelteinflüsse, wie Wind und Helligkeitswert zu reagieren. Dafür wird ein Sensor-Modul gebraucht, um die Steuereinheit auch dezentral platzieren zu können. Das gesamte System sollte möglichst einfach installierbar und anpassbar sein, sowie durch Funkbetrieb Kabelleitungen zwischen den Untersystemen vermeiden und eine einfache Einrichtung gewährleisten.

% Welche Lösung wurde gewählt?
Um ein eben beschriebenes System zu gewährleisten, wurde das System in drei Untersysteme gegliedert. Die Zentralsteuerung wurde mit einem Raspberry Pi 3 realisiert, der mit Raspbian Stretch aufgesetzt wurde und mit HomeAssistant als Smarthome Verwaltung konfiguriert wurde. Um eine schnelle Software Installation zu garantieren, wurde ein Ansible Skript eingesetzt. Zusätzlich wurde Unbound für eine privatere DNS Abfrage installiert. Die Sensorwerte wurden über einen Nordic Mikrocontroller eingelesen und über ein Thread Netzwerk an den Border Router mittels MQTT übermittelt. Dieser konnte dann auf Grund dieser Daten mittels MQTT einen ESP32 Befehle zum Jalousien Stand senden.
\chapter{Anforderungen}
\label{cha:Anforderungen}

% Tabelle Requirements
% UseCase(s)

\section{Requirements}
Die Anforderungen an gesamte System wurden in Requirements geschrieben. Zusätzlich dazu wurde für Jedes eine Kategorie festgelegt, für die es gilt.

\textit{Kategorie:} \textbf{F}(Funktional); \textbf{IF}(Interface);
 \textbf{Q}(Quality)\\
\textit{Verifikationsmethode:} \textbf{S}(Similaris); \textbf{I}(Inspektion); 
\textbf{R}(Review); \textbf{M}(Measurement);
\textbf{A}(Analysis); \textbf{T}(Test)
%Kategorie:
%\begin{itemize}
%	\item F = Funktional
%	\item IF = Interface
%	\item Q = Quality
%\end{itemize}
%Verifikationsmethode:
%\begin{itemize}
%	\item S(imiliaris)
%	\item I(nspection)
%	\item R(eview)
%	\item M(easurement)
%	\item A(nalysis)
%	\item T(est)
%\end{itemize}


% Please add the following required packages to your document preamble:
% \usepackage{booktabs}
% \usepackage[table,xcdraw]{xcolor}
% If you use beamer only pass "xcolor=table" option, i.e. \documentclass[xcolor=table]{beamer}
% \usepackage{longtable}
% Note: It may be necessary to compile the document several times to get a multi-page table to line up properly
\begin{longtable}[ht]{p{0.03\textwidth}  p{0.8\textwidth} p{0.04\textwidth} p{0.04\textwidth}}
	
	\toprule
	\rowcolor[HTML]{FFFC9E} 
	{\color[HTML]{333333} \textbf{Nr.}} & {\color[HTML]{333333} \textbf{Beschreibung}}  & {\color[HTML]{333333} \textbf{Kat.}} & {\color[HTML]{333333} \textbf{VM}} \\* \midrule
	\endfirsthead
	%
	\multicolumn{4}{c}%
	{{\bfseries Table \thetable\ continued from previous page}} \\
	\toprule
	\rowcolor[HTML]{FFFC9E} 
	{\color[HTML]{333333} \textbf{Nr.}} & {\color[HTML]{333333} \textbf{Beschreibung}}  & {\color[HTML]{333333} \textbf{Kat.}} & {\color[HTML]{333333} \textbf{VM}} \\* \midrule
	\endhead
	%
	1 & Das System soll modular in einen Regler, ein Sensormodul und eine Aktorsteuerung aufgeteilt sein                                                                                                                                                                                          & F & T \\* \midrule
	2 & Das System soll die Position einer Jalousie zwischen 0\% und 100\% einstellen können                                                                                                                                          & F& T                                   \\* \midrule
	3 & Der Sollwert für die Jalousieposition soll mit Home Assistant (Smarthome Plattform) manuell einstellbar sein                                                                                                & IF & T                                  \\* \midrule
	4 & Der Sollwert für die Jalousieposition soll eine Schnittstelle zur manuellen Einstellung der Jalousieposition bieten sein                                                                                                                                    & IF& T                                   \\* \midrule
	5 & Der Sollwert für die Jalousieposition soll entsprechend der Helligkeit einstellbar sein                                                                                                                                      & IF& T                                   \\* \midrule
	6 & Der Sollwert für die Jalousieposition soll entsprechend der Windstärke einstellbar sein                                                                                                                                    & IF& T                                   \\* \midrule
	7 & Das System soll aktuelle Sensorwerte ausgeben                                                                                                                                  & IF& T                                   \\* \midrule
	8 & Das System soll die neue Jalousieposition ausgeben                                                                                                                                       & IF& T                                   \\* \midrule
	9 & Die aktuelle Helligkeit soll mit Homeassistant angezeigt werden können                                                                                                                                 & IF& T                                   \\* \midrule
	10 & Die aktuelle Windstärke soll mit Homeassistant angezeigt werden können                                                                                                                                    & IF& T                                   \\* \midrule
	11 & Die neue Jalousieposition soll mit Homeassistant angezeigt werden können 	                                                                                                                                     & IF& T                                   \\* \midrule
	12 & Das Sensormodul soll netzunabhängig betrieben werden.                                                                                                                                  & F& I                                   \\* \midrule
	13 & Die Batterie des Sensormoduls soll mindestens 1 Jahr überdauern                                                                                                                                     & F& T                                   \\* \midrule
	14 & Die Batterie des Sensormoduls soll austauschbar sein                                                                                                                                      & F& I                                  \\* \midrule
	15 & Das System soll 1 mal pro Sekunde aktualisiert werden                                                                                                                                      & F& T                                   \\* \midrule
	16 & Das System (bis auf das Sensormodul) soll über eine externe Spannungsversorgung betrieben werden                                                                                                                                     & F& T                                   \\* \midrule
	17 & Die aktuellen Sensorwerte sollen vom Sensormodul gemessen werden                                                                                                                                      & IF& T                                   \\* \bottomrule
\end{longtable}

\section{UseCase}
Die Abbildung~\ref{fig:UseCase} zeigt den normalen Anwendungsfall des Systems. Der Benutzer kann Einfluss auf die Jalousie Software nehmen, mittels HomeAssistent. Von diesem aus gehen die Befehle an den Regler und es werden aktuelle Informationen an die Statusanzeige gesendet. Der Regler wiederum entscheidet, ob eine Handlung nötig wird und kann dementsprechend Befehle an den Aktor senden, der den Motor der Jalousie steuert. Er kann auch seinen aktuellen Status an die Statusanzeige senden. Außer dem Benutzer kann zusätzlich eine Wartungsperson direkt auf die Jalousie Steuerung zugreifen und änderungen vornehmen.
\begin{figure}[hbt]
	\centering
	\includegraphics[width=1\linewidth]{images/UseCase}
	\caption[UseCase Diagramm]{UseCase Diagramm für den normalen Anwendungsfall.}
	\label{fig:UseCase}
\end{figure}
\chapter{Planung des Projekts}
\label{cha:Planung}

\section{Arbeitspakete}

Es bietet sich an das Projekt in mehrere Arbeitspakete aufzuteilen. Hierzu bietet es sich an sich an den verwendeten Codemodulen zu orientieren. Dies sind in diesem Projekt das Regler Module, Das Actuator Module, das Sensor Module und Codeteile wie der Wifi-Task\nomenclature{Wifi}{Wireless Fidelity} die für mehrere Code Module in gleicher Form benötigt werden. Innerhalb der größeren Module Actuator Module und Sensor Module bietet es sich an eine weiter Unterteilung entsprechend der Tasks vorzunehmen. Die Grafik \ref{fig:arbeitspakete} zeigt die genaue Aufteilung der Arbeitspakete.

\begin{figure}[hbt]
	\centering
	\includegraphics[width=1\linewidth]{images/arbeitspakete}
	\caption[Arbeitspakete]{Arbeitspakete}
	\label{fig:arbeitspakete}
\end{figure}

\section{Zeitplan}
Das Grafik \ref{fig:time_plan} zeigt den Zeitbedarf für die einzelnen Arbeitspakete in Form eines Gant Diagramms. Der Zeitbedarf für die einzelnen Arbeitspakete kann ebenfalls aus der Grafik \ref{fig:arbeitspakete} abgelesen werden.

\begin{figure}[hbt]
	\centering
	\includegraphics[width=1\linewidth]{images/Gant}
	\caption[Zeitplanung als Gant Diagramm]{Zeitplanung als Gant Diagramm}
	\label{fig:time_plan}
\end{figure}


\section{Kostenschätzung}
\label{cha:Planung_cost}
In der Planungsphase wurde angenommen, dass mit anderen Mikrocontrollern und damit einem höheren Aufwand gearbeitet wird. Damit ergeben sich gesamt 2200 SLOC\nomenclature{SLOC}{Source Lines Of Code}. Die Vorlesung Mikrocontroller umfasst 64h und somit ergibt sich im Schnitt 34 lines/h, die geschrieben werden müssen. Die Kostenschätzung nach COCOMOII\_2000.4 ist in Abbildung~\ref{fig:cost_plan} zu sehen.

\begin{figure}[hbt]
	\centering
	\includegraphics[width=1\linewidth]{images/cost_plan_img}
	\caption[Kostenschätzung Planung]{Kostenschätzung bei der Planung.}
	\label{fig:cost_plan}
\end{figure}

% ---- Literaturverzeichnis ----------

\bibliography{literature/literatur1,literature/literatur2}
										% Einbindung mehrerer Verzeichnisse in einem \bibliography Befehl mit Kommata trennen - keine Leerzeichen nach den Kommata!
\bibliographystyle{alphadin}           	%plain: alphabetisch, unsrt: nach Zitat, alphadin: NameJahr


% -----Ausgabe aller Verzeichnisse ---
\setlength{\parskip}{0.5\baselineskip}
\renewcommand{\indexname}{Sachwortverzeichnis}
%\printindex								% Erzeugen des Indexverzeichnises
\addcontentsline{toc}{chapter}{\indexname}
% Version 0.1 vom 31.10.14 - T. Kibler

% Änderungswünsche bitte an T. Kibler melden

% alle Abkürzungen, die in der Arbeit verwendet werden. Nachfolgend sind allgemeine Abkürzungen aufgelistet. Die für die Arbeit spezifischen Abkürzungen sollten innerhalb des Haupttextes aufgeführt werden. Die Alphabetische Sortierung übernimmt Latex.

\nomenclature{Abb.}{Abbildung}
\nomenclature{bzw.}{beziehungsweise}
\nomenclature{DHBW}{Duale Hochschule Baden-Württemberg}
\nomenclature{ebd.}{ebenda}
\nomenclature{et al.}{at alii}
\nomenclature{etc.}{et cetera}
\nomenclature{evtl.}{eventuell}
\nomenclature{f.}{folgende Seite}
\nomenclature{ff.}{fortfolgende Seiten}
\nomenclature{ggf.}{gegebenenfalls}
\nomenclature{Hrsg.}{Herausgeber}
\nomenclature{Tab.}{Tabelle}
\nomenclature{u. a.}{unter anderem}
\nomenclature{usw.}{und so weiter}
\nomenclature{vgl.}{vergleiche}
\nomenclature{z. B.}{zum Beispiel}
\nomenclature{z. T.}{zum Teil}				% Datei mit allgemeinen Abkürzungen laden
\renewcommand{\nomname}{Verzeichnis verwendeter Formelzeichen und Abkürzungen}
\setlength{\nomlabelwidth}{.20\hsize}
\renewcommand{\nomlabel}[1]{#1 \dotfill}
\setlength{\nomitemsep}{-\parsep}
%\printnomenclature						% Erzeugen des Abkürzungsverzeichnises, siehe auch Inhalt der Datei pages/abkuerzungen.tex
\cleardoublepage
%\renewcommand{\glossaryname}{Glossar}
%\printglossaries
%\cleardoublepage
\listoffigures 							% Erzeugen des Abbildungsverzeichnisses 
\cleardoublepage
\listoftables 							% Erzeugen des Tabellenverzeichnisses
\cleardoublepage

% -----Anhang ------------------------

\begin{appendix}
\clearpage
%\pagenumbering{Roman}					% große, römische Seitenzahlen für Anhang, falls gewünscht
\addchap{Anhang im Zip Archiv}
\setcounter{chapter}{1}
\setcounter{section}{0}
\setcounter{table}{0}
\setcounter{figure}{0}

\begin{itemize}
	\item README
	\item Blind Controller (Ordner mit Quellcode)
	\item Sensor Modul (Ordner mit Quellcode)
	\item Aktor Modul (Ordner mit Quellcode)
\end{itemize}
\addchap{Anhang E}
\setcounter{chapter}{5}
\setcounter{section}{0}
\setcounter{table}{0}
\setcounter{figure}{0}

\section{Wichtige \LaTeX -Befehle}

\begin{tabbing}
\hspace*{0cm} \= \hspace{0.25\linewidth} \= \+\kill
\textbackslash \textit{label}\{\}	\> Definition eines Labels, auf welches referenziert werden kann\\ 
	\> z.B.: \textbackslash \textit{label}\{fig:MyImage\}\\ 
\textbackslash \textit{ref}\{\}	\> Setzen einer Referenz zu einem Label\\
\textbackslash \textit{pageref}\{\}	\> Gibt die Seitenzahl zu einer Referenz zurück\\
	\> z.B.: Tabelle\~{}\textbackslash \textit{ref}\{tab:messdaten\} fasst die Messergebnisse zusammen.\\ 
\textbackslash \textit{cite}\{\}	\> Literaturreferenz einfügen\\
\textbackslash \textit{cite}[S. x]\{\}	\> Literaturreferenz mit Angabe einer Seitenzahl \glqq x\grqq~einfügen\\

\textbackslash \textit{footnote}\{\}	\> Fußnote einfügen\\ 
\~{}	\> Einfügen eines geschützten Leerzeichens\\ 
\textdollar \textit{Formel} \textdollar	\> Eingabe einer Formel im Text\\
\textbackslash \textit{nomenclature}\{a.\}\{ab\}	\> Aufnahme der Abkürzung \glqq a.\grqq~für \glqq ab\grqq~in das Abkürzungsverzeichnis.\\
\textbackslash \textit{index}\{Obst!Birne\} \> Aufnahme des Begriffs \glqq Birne\grqq~in den Index unter \glqq Obst\grqq. \index{Obst!Birne} \\
\textbackslash \textit{clearpage}	\> Ausgabe aller Gleitobjekte und Umbruch auf neue Seite\\ 
\end{tabbing}

\clearpage

\section{Vorlagen für \LaTeX Umgebungen}

\subsection{Listen und Aufzählungen}

Es gibt folgende Listentypen. Die wichtigsten:

\begin{itemize}
	\item Einfache Liste mit \textit{itemize}-Umgebung
	\item ...
\end{itemize}

\begin{enumerate}
	\item Nummerierte Liste mit \textit{enumerate}-Umgebung
	\item ...
\end{enumerate}

\begin{enumerate}[label=\alph*.]
	\item wobei man bei der \textit{enumerate}-Umgebung leicht die Art der Numerierung ändern kann,
	\item ...
\end{enumerate}

und durch verschachtelte Umgebungen verschiedene Aufzählungsebenen darstellen kann:

\begin{enumerate}[label=\alph*)]
	\item Erster Aufzählungspunkt der ersten Ebene
	\item ...
	\begin{itemize}
		\item Erster Punkt der zweiten Ebene
		\item Zweiter Punkt der zweiten Ebene
	\end{itemize}
	\item Das sollte an Beispielen zunächst einmal genügen.
\end{enumerate}

\clearpage

\subsection{Bilder und Grafiken}

Bilder können als PDF-, JPG-, und PNG-Bilder in \LaTeX eingebunden werden. Damit eine Grafik in hoher Qualität dargestellt wird, sollte das Dateiformat der Grafik vektorbasiert sein, d.h. als PDF-Datei vorliegen. Viele Zeichenprogramme unterstützen einen PDF-Export (z.B. GIMP, Adobe Illustrator, etc.). Für Grafiken aus PowerPoint sei folgende Vorgehensweise beim Export empfohlen:

\begin{enumerate}
	\item Die gewünschte Grafik in PowerPoint zeichnen.
	\item Gewünschten Bildbereich markieren, rechte Maustaste klicken und \glqq Als Grafik speichern ...\grqq~wählen.
	\item Grafik im Format EMF abspeichern. Das EMF-Format ist vektorbasiert.\footnote{Mit dem Mac kann in PowerPoint die Grafik direkt im PDF-Format exportiert werden. Die weiteren Schritte entfallen daher.}
	\item Mit dem Programm XnView die Grafik im EMF-Format in PDF wandeln und abspeichern.
	\item Die so erzeugte PDF-Datei enthält eine vektorbasierte Grafik und kann in \LaTeX eingebunden werden.
\end{enumerate}

Abbildung~\ref{fig:MyImage} zeigt ein Beispielbild einer Grafik, welche aus PowerPoint exportiert wurde.

\begin{figure}[hbt]
	\centering
	\includegraphics[width=0.3\linewidth]{images/MyImage}
	\caption[Beispiel für die Einbindung eines Bildes.]{Beispiel für die Einbindung eines Bildes (PDF-, JPG-, und PNG-Bilder können eingebunden werden).}
	\label{fig:MyImage}
\end{figure}

Der Quellcode des Beispielbildes aus Abbildung~\ref{fig:MyImage} ist in Listing~\ref{lst:fig} zu sehen.

\begin{lstlisting}[caption=Quellcode der Abbildung~\ref{fig:MyImage}.,label=lst:fig]
\begin{figure}[hbt]				% here, bottom, top
\centering						% Zentrierung
\includegraphics[width=0.6\linewidth]{images/MyImage}		
\caption[Beispiel für die Einbindung eines Bildes.]{Beispiel für die Einbindung eines Bildes (PDF-, JPG-, und PNG-Bilder können eingebunden werden).}
\label{fig:MyImage}
\end{figure}
\end{lstlisting}

\clearpage

Grafiken mit Tikz werden mit dem \textit{input}-Befehl in die \textit{figure}-Umgebung geladen:

\begin{figure}[hbt]
	\centering
	\begin{tikzpicture}
		\begin{axis}[scale=1.3,legend entries={Messwerte mit Fehlerbalken,
			$\pgfmathprintnumber{\pgfplotstableregressiona} \cdot x  
			\pgfmathprintnumber[print sign]{\pgfplotstableregressionb}$}, legend style={draw=none},legend style={at={(0.01,0.98)},anchor=north west},xlabel=Stromstärke $I \; \mathrm{ \lbrack mA \rbrack}$,ylabel=Spannung $U \; \mathrm{ \lbrack V \rbrack}$]
		\addlegendimage{mark=*,blue}
		\addlegendimage{no markers,red}
\addplot+[error bars/.cd, y dir=both,y explicit]
table[x=x,y=y,y error=errory] 
{pgfplot/messdaten_mitfehler.dat};
\addplot table[mark=none,y={create col/linear regression={y=y}}]
{pgfplot/messdaten_mitfehler.dat};
	\end{axis}
\end{tikzpicture}
	\caption[Diagramm, erstellt mit dem \textit{pgfplot}-Befehlssatz.]{Ein Diagramm, erstellt in der \textit{tikzpicture}-Umgebung mit dem \textit{pgfplot}-Befehlssatz. Das Diagramm stellt Messdaten, deren Fehlerbalken und eine Regressionskurve dar. Die Messdaten werden von einer separaten Datei eingelesen und die Regressionskurve wurde mit \textit{pgfplot} berechnet und erstellt.}
	\label{fig:pgfplot}
\end{figure}

Auch hierzu der Quellcode in Listing~\ref{lst:pgfplot}, in Listing~\ref{lst:tikz} ist der Quellcode der Datei \textit{mess\_fehlerbalken.tex} dargestellt.

\begin{lstlisting}[caption=Quellcode der Abbildung~\ref{fig:pgfplot}.,label=lst:pgfplot]
\begin{figure}[hbt]
\centering
\begin{tikzpicture}
		\begin{axis}[scale=1.3,legend entries={Messwerte mit Fehlerbalken,
			$\pgfmathprintnumber{\pgfplotstableregressiona} \cdot x  
			\pgfmathprintnumber[print sign]{\pgfplotstableregressionb}$}, legend style={draw=none},legend style={at={(0.01,0.98)},anchor=north west},xlabel=Stromstärke $I \; \mathrm{ \lbrack mA \rbrack}$,ylabel=Spannung $U \; \mathrm{ \lbrack V \rbrack}$]
		\addlegendimage{mark=*,blue}
		\addlegendimage{no markers,red}
\addplot+[error bars/.cd, y dir=both,y explicit]
table[x=x,y=y,y error=errory] 
{pgfplot/messdaten_mitfehler.dat};
\addplot table[mark=none,y={create col/linear regression={y=y}}]
{pgfplot/messdaten_mitfehler.dat};
	\end{axis}
\end{tikzpicture}
\caption[Diagramm, erstellt mit dem \textit{pgfplot}-Befehlssatz.]{Ein Diagramm, erstellt in der \textit{tikzpicture}-Umgebung mit dem \textit{pgfplot}-Befehlssatz. Das Diagramm stellt Messdaten, deren Fehlerbalken und eine Regressionskurve dar. Die Messdaten werden von einer separaten Datei eingelesen und die Regressionskurve wurde mit \textit{pgfplot} berechnet und erstellt.}
\label{fig:pgfplot}
\end{figure}
\end{lstlisting}

\begin{lstlisting}[caption=Quellcode der Datei \textit{mess\_fehlerbalken.tex}.,label=lst:tikz]
\begin{tikzpicture}
\begin{axis}[scale=1.3,legend entries={Messwerte mit Fehlerbalken,
$\pgfmathprintnumber{\pgfplotstableregressiona} \cdot x  
\pgfmathprintnumber[print sign]{\pgfplotstableregressionb}$}, legend style={draw=none},legend style={at={(0.01,0.98)},anchor=north west},xlabel=Stromstärke $I \; \mathrm{ \lbrack mA \rbrack}$,ylabel=Spannung $U \; \mathrm{ \lbrack V \rbrack}$]
\addlegendimage{mark=*,blue}
\addlegendimage{no markers,red}
\addplot+[error bars/.cd, y dir=both,y explicit]
table[x=x,y=y,y error=errory] 
{pgfplot/messdaten_mitfehler.dat};
\addplot table[mark=none,y={create col/linear regression={y=y}}]
{pgfplot/messdaten_mitfehler.dat};
\end{axis}
\end{tikzpicture}
\end{lstlisting}

\clearpage

\subsection{Tabellen}

\begin{table}[hbt]	
	\centering
	\renewcommand{\arraystretch}{1.5}	% Skaliert die Zeilenhöhe der Tabelle
	\captionabove[Liste der verwendeten Messgeräte]{Liste der verwendeten Messgeräte. Die Genauigkeitsangaben beziehen sich auf die Standardabweichung $1\cdot \sigma$.}
	\label{tab:bsp}
	\begin{tabular}{ccccc}
		\textbf{Messgerät} & \textbf{Hersteller} & \textbf{Typ} & \textbf{Verwendung} & \textbf{Genauigkeit}\\ 
		\hline 
		\hline 
		\parbox[t]{0.2\linewidth}{\centering Spannungs-\\versorgung} & Voltmaker & HV2000 & \parbox[t]{0.2\linewidth}{\centering Spannungs-\\versorgung der\\Platine} & $\Delta U = \pm 5 $~mV \\ % Der parbox-Befehl ist erforderlich, damit ein Zeilenumbruch erzeugt werden kann. c-Spalten (zentriert) erlauben nicht automatisch einen Zeilenumpruch. Linksbündig gesetzte p-Spalten erlauben automatisch den Zeilenumbruch.
		Strommessgerät & Currentcount & Hotamp 16 & \parbox[t]{0.2\linewidth}{ \centering Strommessung\\am Versorgungspin\\des µC} & $\Delta I = \pm 0.1$~A \\ 
		\hline 
	\end{tabular} 
\end{table}

Der Quellcode der Beispieltabelle~\ref{tab:bsp} ist in Listing~\ref{lst:tab} zu sehen.

\begin{lstlisting}[caption=Quellcode der Tabelle~\ref{tab:bsp}.,label=lst:tab]
\begin{table}[hbt]	
\centering
\renewcommand{\arraystretch}{1.5}	% Skaliert die Zeilenhöhe der Tabelle
\captionabove[Liste der verwendeten Messgeräte]{Liste der verwendeten Messgeräte. Die Genauigkeitsangaben beziehen sich auf die Standardabweichung $1\cdot \sigma$.}
\label{tab:bsp}
\begin{tabular}{ccccc}
\textbf{Messgerät} & \textbf{Hersteller} & \textbf{Typ} & \textbf{Verwendung} & \textbf{Genauigkeit}\\ 
\hline 
\hline 
\parbox[t]{0.2\linewidth}{\centering Spannungs-\\versorgung} & Voltmaker & HV2000 & \parbox[t]{0.2\linewidth}{\centering Spannungs-\\versorgung der\\Platine} & $\Delta U = \pm 5 $~mV \\ % Der parbox-Befehl ist erforderlich, damit ein Zeilenumbruch erzeugt werden kann. c-Spalten (zentriert) erlauben nicht automatisch einen Zeilenumpruch. Linksbündig gesetzte p-Spalten erlauben automatisch den Zeilenumbruch.
Strommessgerät & Currentcount & Hotamp 16 & \parbox[t]{0.2\linewidth}{ \centering Strommessung\\ am Versorgungspin\\ des \textmu C} & $\Delta I = \pm 0.1$~A \\ 
\hline 
\end{tabular} 
\end{table}
\end{lstlisting}

\clearpage

\subsection{Formeln}

Formeln lassen sich in \LaTeX~ganz einfach schreiben. Es gibt unterschiedliche Umgebungen zum Schreiben von Formeln. Z.B. direkt im Text $v=s/t$ oder abgesetzt

\[F=m \cdot a\]

oder auch, wie in wissenschaftlichen Dokumenten üblich, nummeriert

\begin{equation}
P=\frac{U^2}{R} \quad .
\label{eqn:leistung}
\end{equation}

Mit einem Label in Formel~\ref{eqn:leistung} lassen sich natürlich auch Formeln im Text referenzieren. \LaTeX~verwendet im Formelmodus einen eigenen Schriftsatz, welcher entsprechend der gängigen Konventionen kursive Zeichen verwendet. Sollen im Formelmodus Einheiten in normaler Schriftart eingefügt werden, dann kann dies über den Befehl \textbackslash \textit{mathrm}\{\} erwirkt werden, wie im Quellcode von Formel~\ref{eqn:leistungMitEinh} zu sehen ist.

\begin{equation}
P=\frac{U^2}{R} = \frac{\left( 100~\mathrm{V}\right)^2}{100~\Omega} = 100~\mathrm{W}\quad .
\label{eqn:leistungMitEinh}
\end{equation}

Zum direkten Vergleich sind die Einheiten in Formel~\ref{eqn:leistungMitEinhfalsch} falsch dargestellt:

\begin{equation}
P=\frac{U^2}{R} = \frac{\left( 100~V\right)^2}{100\,\varOmega} = 100\,W
\label{eqn:leistungMitEinhfalsch}
\end{equation}

Das sind nur ein paar wenige Beispiele und es gibt sehr viele Packages, um Besonderheiten in Formeln realisieren zu können, z.B. mehrzeilige Formeln mit vertikaler Ausrichtung. Nennen Sie Formeln nur, wenn diese zum besseren Verständnis auch wirklich nützlich sind.

Folgende Befehle sind innerhalb von Formel-Umgebungen nützlich:
\begin{tabbing}
	\hspace*{0cm} \= \hspace{0.35\linewidth} \= \+\kill
	\textbackslash \textit{text}\{\}	\> Damit kann in Formel-Umgebung Text geschrieben werden.\\ 
	\textbackslash, \textbackslash: \textbackslash; oder \textbackslash quad und \textbackslash qquad \> Zusätzlichen Abstand zwischen Symbolen einfügen.\\
	\textbackslash \textit{notag} \> Nummerierung einer bestimmten Formel ausschalten.
\end{tabbing}

Abschließend nochmals ein kleines Beispiel:

\begin{eqnarray}
\sum\limits_{n=1}^\infty f\left(x_n\right)\cdot \Delta x=  \lim\limits_{\Delta x \rightarrow 0} \frac{f\left(x_0+\Delta x\right)-f\left(x_0\right)}{\Delta x} = \frac{\diff f}{\diff x} = \dot{f}(x)
\end{eqnarray}		% Zeile auskommentieren bei finalem Dokument!
\end{appendix}


\end{document}